% paper-annotated.tex -- REVIEW-ONLY annotated build of the MDPI paper
% Identical content to paper.tex; adds cosmetic per-reviewer highlights and
% Word-style margin comments (Reviewer 1 = orange, Reviewer 2 = blue) that map
% each review comment to the passage resolving it.
%
% The clean manuscript (paper.tex + sections/*.tex) is unchanged. This build
% inputs annotated copies from sections-annotated/.
% For a markup-free build, edit reviewer-annotation-spec.tex to load
% \usepackage[final]{changes}.
%
% Compiled with: pdflatex -> bibtex -> pdflatex -> pdflatex

\documentclass[symmetry,article,submit,pdftex,moreauthors]{Definitions/mdpi}

% =================================================================
% MDPI internal commands
\firstpage{1}
\makeatletter
\setcounter{page}{\@firstpage}
\makeatother
\pubvolume{1}
\issuenum{1}
\articlenumber{0}
\pubyear{2026}
\copyrightyear{2026}

% =================================================================
% Additional packages not in mdpi.cls
\usepackage{pgfplots}
\usetikzlibrary{shapes.geometric, shapes.misc, arrows.meta, positioning, fit,
                decorations.pathreplacing, calc, backgrounds,
                matrix, chains, scopes}
\usepgfplotslibrary{groupplots}
\pgfplotsset{compat=1.18}
\usepackage{adjustbox}

% Code listings
\lstset{
  basicstyle=\small\ttfamily,
  keywordstyle=\bfseries,
  commentstyle=\itshape,
  frame=single,
  breaklines=true,
  captionpos=b,
  numbers=left,
  numberstyle=\tiny,
  language=Python
}

% Theorem environments (defined in mdpi.cls, but custom ones needed)
\newtheorem{definition}{Definition}
\newtheorem{proposition}{Proposition}
\newtheorem{theorem}{Theorem}

% =================================================================
% Custom commands
\newcommand{\EFE}{\ensuremath{\mathcal{G}}}
\newcommand{\FE}{\ensuremath{\mathcal{F}}}
\newcommand{\KL}[2]{\ensuremath{D_{\mathrm{KL}}\!\left(#1 \,\|\, #2\right)}}
\newcommand{\E}[1]{\ensuremath{\mathbb{E}\!\left[#1\right]}}
\newcommand{\softmax}{\ensuremath{\operatorname{softmax}}}
\newcommand{\Attn}{\ensuremath{\operatorname{Attn}}}
\newcommand{\MLP}{\ensuremath{\operatorname{MLP}}}
\newcommand{\argmin}{\ensuremath{\operatorname*{arg\,min}}}
\newcommand{\abs}[1]{\ensuremath{\left|#1\right|}}
\newcommand{\norm}[1]{\ensuremath{\left\|#1\right\|}}
\newcommand{\vect}[1]{\ensuremath{\bm{#1}}}
\newcommand{\mat}[1]{\ensuremath{\mathbf{#1}}}
\newcommand{\ACDC}{ACDC}
\newcommand{\EAP}{EAP}
\newcommand{\ACD}{ACD}
\newcommand{\pymdp}{\texttt{pymdp}}

% =================================================================
% Reviewer annotation markup (review-only; see docs/overleaf-comments.md)
% ============================================================
% reviewer-annotation-spec.tex
% Cosmetic, review-only markup: highlights reviewer comments and their
% resolutions directly on the manuscript as Word-style margin notes.
% Reviewer 1 = orange, Reviewer 2 = blue.
%
% This file is \input only by paper-annotated.tex. The clean manuscript
% (paper.tex + sections/*.tex) does not load it and is unchanged.
%
% To produce a clean (markup-free) build, load changes with [final]:
%   \usepackage[final]{changes}
% ============================================================

% Float-free margin notes (Word-style side comments). marginnote is used
% instead of todonotes so that comments next to MDPI's [H] tables/figures
% are never dropped ("Float(s) lost").
\usepackage{marginnote}
\renewcommand*{\marginfont}{\scriptsize\raggedright}

% Per-author change tracking with margin comments
\usepackage[
  markup=underlined,
  authormarkup=superscript,
  commentmarkup=default
]{changes}

% Reviewers as "authors" with distinct colours
\definechangesauthor[name={Reviewer 1}, color=orange]{Rone}
\definechangesauthor[name={Reviewer 2}, color=blue]{Rtwo}

% Math/macro-safe markup: highlighted spans only recolour the text (the
% per-author colour is applied automatically), so they may safely contain
% inline math ($...$), \emph, \texttt, \ref, etc. without soul/ulem breakage.
% "authorcolor" is set by changes to the current reviewer's colour before the
% markup expands; wrapping in \textcolor keeps per-reviewer colour while being
% safe for inline math and macros inside the span.
\sethighlightmarkup{\textcolor{authorcolor}{#1}}
\setaddedmarkup{\textcolor{authorcolor}{#1}}

% Route reviewer comments to float-free margin notes (coloured per author).
\setcommentmarkup{\marginnote{\textcolor{authorcolor}{[#1]}}}

% ===== Convenience macros (see docs/overleaf-comments.md) =====
% <comment> is a short "Cxx: comment -> resolution" string shown in the margin.

% --- Reviewer 1 (orange) ---
\newcommand{\RoneAdd}[2][]{\added[id=Rone, comment={#1}]{#2}}
\newcommand{\RoneDel}[2][]{\deleted[id=Rone, comment={#1}]{#2}}
\newcommand{\RoneRep}[3][]{\replaced[id=Rone, comment={#1}]{#3}{#2}}
\newcommand{\RoneHl}[2][]{\highlight[id=Rone, comment={#1}]{#2}}
\newcommand{\RoneCom}[1]{\comment[id=Rone]{#1}}

% --- Reviewer 2 (blue) ---
\newcommand{\RtwoAdd}[2][]{\added[id=Rtwo, comment={#1}]{#2}}
\newcommand{\RtwoDel}[2][]{\deleted[id=Rtwo, comment={#1}]{#2}}
\newcommand{\RtwoRep}[3][]{\replaced[id=Rtwo, comment={#1}]{#3}{#2}}
\newcommand{\RtwoHl}[2][]{\highlight[id=Rtwo, comment={#1}]{#2}}
\newcommand{\RtwoCom}[1]{\comment[id=Rtwo]{#1}}


% =================================================================
% MDPI metadata (before \begin{document})

\Title{Active Circuit Discovery: A Multi-Action POMDP Agent
       for Causal Feature Identification\\in Transformer Attribution Graphs}

\newcommand{\orcidauthorA}{0000-0000-0000-000X}

\Author{Sharath Sathish $^{1}$\orcidA{}, Mominul Ahsan $^{2}$\orcidA{} and Majid Latifi $^{2,}$*\orcidA{}}

\AuthorNames{Sharath Sathish, Mominul Ahsan and Majid Latifi}

\address{%
$^{1}$ \quad BP International Limited, London TW16 7BP, UK; sharath.sathish@gmail.com \\
$^{2}$ \quad Department of Computer Science, University of York, York YO10 5GH, UK; mominul.ahsan2@gmail.com; majid.latifi@york.ac.uk}

\corres{majid.latifi@york.ac.uk}

\abstract{\input{sections-annotated/abstract}}

\keyword{mechanistic interpretability; active inference; circuit discovery; transcoders; Expected Free Energy; POMDP}

% =================================================================
\begin{document}

\maketitle

% =================================================================
\section{Introduction}
\label{sec:intro}
\input{sections-annotated/introduction}

% =================================================================
\section{Background and Related Work}
\label{sec:background}
\input{sections-annotated/background}

% =================================================================
\section{The Active Circuit Discovery Framework}
\label{sec:method}
\input{sections-annotated/methodology}

% =================================================================
\section{Experimental Setup}
\label{sec:setup}
\input{sections-annotated/experimental_setup}

% =================================================================
\section{Results}
\label{sec:results}
\input{sections-annotated/results}

% =================================================================
\section{Discussion}
\label{sec:discussion}
\input{sections-annotated/discussion}

% =================================================================
\section{Conclusions and Future Work}
\label{sec:conclusion}
\input{sections-annotated/conclusion}

% =================================================================
% MDPI mandatory end sections
\vspace{6pt}

\authorcontributions{Conceptualization, S.S.; methodology, S.S.; software, S.S.;
validation, S.S., M.A.\ and M.L.; formal analysis, S.S.; investigation, S.S.;
resources, S.S.; data curation, S.S.; writing---original draft preparation, S.S.;
writing---review and editing, S.S., M.A.\ and M.L.; visualization, S.S.;
supervision, M.A.\ and M.L.; project administration, S.S.
All authors have read and agreed to the published version of the manuscript.}

\funding{This research received no external funding.}

\institutionalreview{Not applicable.}

\informedconsent{Not applicable.}

\dataavailability{Code, experiment notebooks, and raw JSON results are
publicly available at \url{https://github.com/SharathSPhD/ActiveCIrcuitDiscovery}.
The Gemma-2-2B model is available under the Gemma licence from
\url{https://huggingface.co/google/gemma-2-2b}.
Llama-3.2-1B is available under the Meta Llama licence from HuggingFace Hub
(\url{https://huggingface.co/meta-llama/Llama-3.2-1B}).}

\acknowledgments{The authors thank the Anthropic interpretability team
for releasing the \texttt{circuit-tracer} library and GemmaScope
transcoder weights that made this work possible. The authors also
gratefully acknowledge the University of York and BP for providing the
academic and institutional environment that enabled them to undertake
this research.

The authors used various AI tools to enhance the language and
readability of the paper during the writing process. After utilizing
these tools, the authors evaluated and performed any necessary editing
of the text. The authors are solely responsible for the fundamental
study, results, and research findings.}

\conflictsofinterest{The authors declare no conflicts of interest.}

% =================================================================
\abbreviations{Abbreviations}{
The following abbreviations are used in this manuscript:\\
\noindent
\begin{tabular}{@{}ll}
ACD   & Active Circuit Discovery \\
ACDC  & Automated Circuit Discovery \\
EAP   & Edge Attribution Patching \\
EFE   & Expected Free Energy \\
FEP   & Free Energy Principle \\
IOI   & Indirect Object Identification \\
KL    & Kullback--Leibler (divergence) \\
LLM   & Large Language Model \\
MLP   & Multilayer Perceptron \\
POMDP & Partially Observable Markov Decision Process \\
SAE   & Sparse Autoencoder \\
\end{tabular}
}

% =================================================================
% Appendices
\appendixtitles{yes}
\appendixstart
\appendix
\input{sections-annotated/appendix}

% =================================================================
\reftitle{References}

\begin{adjustwidth}{0cm}{0cm}

\bibliographystyle{mdpi}
\bibliography{references}

\end{adjustwidth}

\PublishersNote{}

\end{document}
